\chapter{Gestione dei dati, persistenza ed esportazione}
\label{chap:dati-export}

\section{Panoramica delle sorgenti dati}

Il progetto supporta quattro sorgenti dati principali, ma con accessi differenti a
seconda dell'interfaccia usata. La GUI puo' lavorare con \emph{PlayTennis}, con il
dataset \emph{Iris} incluso nelle risorse, con file CSV esterni e con tabelle MySQL. Il
client testuale, invece, accede direttamente soltanto al database e ai file
\module{.dmp}. Questa asimmetria non e' accidentale: la GUI e' stata pensata come
front-end completo per l'uso locale, mentre il client testuale e' un front-end minimo per
il backend remoto.

\section{Caricamento da CSV}

\subsection{Formato atteso}

Il parser CSV implementato in \module{Data} richiede una struttura semplice. La prima
riga deve contenere i nomi degli attributi, separati da virgole; le righe successive
devono contenere esattamente lo stesso numero di campi. Le righe vuote vengono ignorate,
mentre l'assenza di dati o l'incoerenza nel numero di colonne genera un'eccezione.

Il parser usa una separazione testuale basata su \module{split(",")}. Questo dettaglio ha
un effetto importante: il formato supportato e' un CSV semplice, senza gestione avanzata
di campi quotati contenenti virgole interne. Per un uso affidabile del sistema e' quindi
opportuno fornire file tabulari puliti, con valori atomici e senza dipendere da
convenzioni CSV piu' sofisticate.

\subsection{Inferenza dei tipi}

La classificazione degli attributi avviene automaticamente per colonna. Un attributo
viene considerato continuo solo se tutti i valori non mancanti sono numerici e i valori
distinti superano la soglia di cinque. In caso contrario l'attributo e' trattato come
discreto. I marker \module{?}, stringa vuota e \module{NA} vengono ignorati durante la
sola fase di inferenza.

Questo comportamento va interpretato con cautela. Il fatto che un valore mancante venga
ignorato nell'inferenza non significa che il sistema implementi una vera strategia di
gestione dei missing values. Se una colonna viene classificata come continua, valori non
numerici presenti nelle righe possono comunque causare errori in fasi successive, quando
la tupla viene convertita negli oggetti usati per il clustering. In termini pratici, il
manuale raccomanda di usare CSV privi di valori mancanti o di normalizzarli prima del caricamento.

\section{Dataset standard e dati hardcoded}

Il dataset \emph{PlayTennis} e' incorporato direttamente nel costruttore di default della
classe \module{Data}. Si tratta di un dataset piccolo, utile per test rapidi e per
comprendere la struttura di un clustering su attributi prevalentemente discreti. Il
dataset \emph{Iris}, invece, e' caricato da una risorsa CSV interna alla GUI tramite
\module{DatasetLoader}. Per renderlo compatibile con il parser esistente, il loader copia
temporaneamente il file delle risorse su disco, crea l'oggetto \module{Data} e poi
rimuove il file temporaneo.

Questi due dataset svolgono funzioni diverse. \emph{PlayTennis} e' ideale per una
verifica rapida del comportamento dell'algoritmo e per testare l'infrastruttura.
\emph{Iris} e' invece piu' adatto alla visualizzazione grafica, perche' contiene
attributi continui ben interpretabili sugli assi del grafico 2D.

\section{Caricamento da database}

\subsection{Schema relazionale supportato}

Quando i dati provengono da MySQL, il backend non copia indiscriminatamente tutte le
colonne della tabella. \module{TableSchema} ispeziona i metadati JDBC e ignora le colonne
auto-increment, tipicamente usate come chiavi primarie tecniche. Le colonne numeriche
vengono trattate come candidate ad attributi continui, mentre quelle testuali vengono
interpretate come discrete.

Il file \filepath{setup_database.sql} fornisce un esempio concreto di tabelle
compatibili. In tutti i casi la colonna \module{id} ha solo funzione identificativa e non
entra nel dataset di clustering. L'utente deve quindi ragionare sulla propria tabella
come su una matrice di attributi descrittivi, non come su una normale tabella
amministrativa con chiavi e metadati ausiliari mescolati alle variabili di interesse.

\subsection{Righe distinte e aggregati}

Il caricamento effettivo viene eseguito da \module{TableData}, che usa una query
\module{SELECT DISTINCT}. Questo significa che tuple duplicate a livello di attributi
descrittivi vengono collassate in una sola occorrenza nel dataset caricato. Inoltre, per
gli attributi numerici vengono calcolati minimo e massimo, indispensabili alla
normalizzazione adottata dal backend.

Per l'utente questo ha due implicazioni. La prima e' che la cardinalita' del dataset
letto dal sistema potrebbe non coincidere con il numero totale di righe della tabella SQL
se esistono duplicati perfetti. La seconda e' che la scala dei valori continui dipende
integralmente dai valori presenti nella tabella al momento del caricamento.

\section{Persistenza in formato \module{.dmp}}

Il formato \module{.dmp} e' la forma di persistenza nativa del progetto. Attraverso il
metodo \module{saveComplete()} del backend vengono serializzati insieme il
\module{ClusterSet}, il dataset originario e il \emph{radius} usato per l'esecuzione.
Questo consente non solo di conservare il risultato finale, ma anche di riaprirlo
successivamente nella GUI con tutte le informazioni necessarie a popolare la schermata
dei risultati.

Il backend conserva comunque compatibilita' con un formato legacy in cui nel file era
presente solo l'insieme dei cluster. In quel caso il caricamento e' ancora possibile, ma
l'applicazione non puo' ricostruire i dettagli completi del dataset. La GUI, infatti,
mostra un avviso quando apre un file di questo tipo e chiarisce che la consultazione
piena dei risultati richiede il formato completo piu' recente.

\section{Esportazione dei risultati}

\subsection{CSV}

L'export CSV della GUI genera un file orientato alla consultazione tabellare.
L'intestazione contiene sempre le colonne \module{ClusterID}, \module{TupleID} e
\module{DistanceFromCentroid}, seguite dai nomi di tutti gli attributi del dataset. Ogni
riga descrive una tupla del clustering finale e riporta la sua distanza dal centroide del
cluster di appartenenza. Questo formato e' particolarmente utile per analisi successive
in fogli elettronici o strumenti esterni.

\subsection{Report TXT}

Il report testuale organizza l'output in due parti principali. La prima raccoglie i
metadati dell'esecuzione: timestamp, radius, tempo di esecuzione e statistiche globali
sui cluster. La seconda passa in rassegna ciascun cluster, riportando dimensione,
centroide, distanze minima/massima/media e una lista delle tuple appartenenti al cluster.
Se il cluster e' molto grande, il report limita la visualizzazione dettagliata alle prime
tuple e segnala quante ne restano escluse dalla stampa.

\subsection{ZIP}

L'archivio ZIP e' pensato come pacchetto completo di consegna. La GUI crea file
temporanei per \module{.dmp}, CSV e report TXT, li inserisce nello ZIP e aggiunge un
\module{README.txt} descrittivo. I nomi dei file interni sono costruiti a partire da un
timestamp, cosi' da rendere evidente il legame tra i diversi artefatti prodotti nella
stessa esportazione.

\subsection{PNG}

L'export PNG riguarda la sola vista grafica dei cluster. Come illustrato nel
Capitolo~\ref{chap:gui}, dedicato all'interfaccia grafica, l'utente puo' salvare il
grafico in formato standard o alta definizione. Questo tipo di esportazione e' distinto
dai tre precedenti: non descrive il clustering in forma tabellare o serializzata, ma ne
offre una rappresentazione visuale.

\section{Indicazioni operative}

Dal punto di vista pratico, la scelta tra i diversi formati dipende dallo scopo. Il file
\module{.dmp} e' la soluzione adatta quando si vuole riaprire in futuro lo stesso
clustering dentro l'applicazione. Il CSV e' preferibile quando si intende lavorare sui
dati con strumenti esterni. Il report TXT e' utile per una lettura umana immediata,
mentre lo ZIP e' indicato quando si desidera archiviare o consegnare un pacchetto
autosufficiente che contenga contemporaneamente persistenza, dati tabellari e
documentazione sintetica dell'esecuzione.
